\documentclass[10pt, aspectratio=169]{beamer}
% Preámbulo modular para presentaciones Beamer (Modelación Estadística)
% Diseño y paleta institucional del Tecnológico de Monterrey

\documentclass[aspectratio=169,xcolor={dvipsnames,table}]{beamer}

\usepackage[utf8]{inputenc}
\usepackage[spanish,mexico]{babel}
\usepackage{amsmath,amssymb,amsfonts,mathtools}
\usepackage{graphicx}
\usepackage{listings}
\usepackage{booktabs}
\usepackage{tikz}
\usetikzlibrary{matrix,arrows,decorations.pathmorphing,shapes.geometric,positioning}

% Paleta Institucional Tecnológico de Monterrey
\definecolor{TecAzul}{HTML}{0039A6}       % Azul TEC: color primario institucional
\definecolor{TecAzulOscuro}{HTML}{1A2E51} % Azul marino oscuro
\definecolor{TecRojo}{HTML}{EC2661}       % Rojo de acento (paleta miTec)
\definecolor{TecGris}{HTML}{646464}       % Gris neutro para comentarios y subtítulos
\definecolor{TecFondoBloque}{HTML}{F4F6F9}% Fondo suave para bloques

% Configuración de Tema Beamer
\usetheme{Boadilla}
\usecolortheme[named=TecAzulOscuro]{structure}
\setbeamercolor{alerted text}{fg=TecRojo}
\setbeamercolor{palette primary}{bg=TecAzulOscuro,fg=white}
\setbeamercolor{palette secondary}{bg=TecAzul,fg=white}
\setbeamercolor{palette tertiary}{bg=TecAzulOscuro!80!black,fg=white}
\setbeamercolor{title}{fg=TecAzulOscuro,bg=TecFondoBloque}
\setbeamercolor{frametitle}{fg=TecAzulOscuro,bg=TecFondoBloque}

% Bloques personalizados elegantes
\setbeamercolor{block title}{bg=TecAzulOscuro,fg=white}
\setbeamercolor{block body}{bg=TecFondoBloque,fg=black}
\setbeamercolor{block title alerted}{bg=TecRojo,fg=white}
\setbeamercolor{block body alerted}{bg=TecRojo!8,fg=black}
\setbeamercolor{block title example}{bg=TecAzul,fg=white}
\setbeamercolor{block body example}{bg=TecAzul!8,fg=black}

% Redondear bordes y añadir sombra sutil
\setbeamertemplate{blocks}[rounded][shadow=true]
\setbeamertemplate{navigation symbols}{}

% Pie de página limpio con número de diapositiva
\setbeamertemplate{footline}{
  \leavevmode%
  \hbox{%
  \begin{beamercolorbox}[wd=.7\paperwidth,ht=2.25ex,dp=1ex,left]{author in head/foot}%
    \hspace*{2ex}\insertshortauthor~~(\insertshortinstitute) \hfill \insertshorttitle
  \end{beamercolorbox}%
  \begin{beamercolorbox}[wd=.3\paperwidth,ht=2.25ex,dp=1ex,right]{date in head/foot}%
    \insertshortdate{}\hspace*{2em}
    \insertframenumber{} / \inserttotalframenumber\hspace*{2ex}
  \end{beamercolorbox}}%
  \vskip0pt%
}

% Configuración de Listings de Python para visualización pedagógica de código
\lstset{
    language=Python,
    basicstyle=\ttfamily\footnotesize,
    keywordstyle=\color{TecAzulOscuro}\bfseries,
    stringstyle=\color{TecRojo},
    commentstyle=\color{TecGris}\itshape,
    showstringspaces=false,
    breaklines=true,
    frame=lines,
    rulecolor=\color{TecAzulOscuro!30},
    backgroundcolor=\color{TecFondoBloque},
    aboveskip=0.5em,
    belowskip=0.5em,
    literate={á}{{\'a}}1 {é}{{\'e}}1 {í}{{\'i}}1 {ó}{{\'o}}1 {ú}{{\'u}}1
             {Á}{{\'A}}1 {É}{{\'E}}1 {Í}{{\'I}}1 {Ó}{{\'O}}1 {Ú}{{\'U}}1
             {ñ}{{\~n}}1 {Ñ}{{\~N}}1 {¿}{{?`}}1 {¡}{{!`}}1
}

% Metadatos institucionales por defecto
\title[Modelación Estadística]{Modelación Estadística para Ciencia de Datos e Ingeniería}
\author[J. Castillo Colmenares]{Juliho David Castillo Colmenares}
\institute[Tec de Monterrey]{Tecnológico de Monterrey \\ Escuela de Ingeniería y Ciencias}
\date{\today}

% Comandos y macros estadísticas compartidas para las presentaciones Beamer (Metropolis)

% Conjuntos numéricos básicos
\newcommand{\R}{\mathbb{R}}
\newcommand{\N}{\mathbb{N}}
\newcommand{\Q}{\mathbb{Q}}
\newcommand{\Z}{\mathbb{Z}}
\newcommand{\set}[1]{\left\{ #1 \right\}}
\newcommand{\abs}[1]{\left|#1\right|}
\newcommand{\norm}[1]{\left\| #1 \right\|}

% Comandos estadísticos y probabilísticos del curso
\newcommand{\Var}{\operatorname{Var}}
\newcommand{\cov}{\operatorname{Cov}}
\newcommand{\corr}{\rho}
\newcommand{\comb}[2]{\begin{pmatrix} #1 \\ #2 \end{pmatrix}}
\newcommand{\s}{\sigma}
\newcommand{\particion}{\mathfrak{P}}
\newcommand{\card}[1]{\#\left( #1 \right)}
\newcommand{\seq}[3]{\set{#1}_{#2}^{#3}}
\newcommand{\E}{\mathbb{E}}
\newcommand{\Prob}{\mathbb{P}}

% Entornos pedagógicos y cajas en español académico natural
\newenvironment{discusion}{%
  \begin{alertblock}{Pregunta de Discusión}%
}{%
  \end{alertblock}%
}

\newenvironment{reflexion}{%
  \begin{alertblock}{Reflexión para el Aula}%
}{%
  \end{alertblock}%
}

\newenvironment{paradoja}{%
  \begin{alertblock}{Paradoja de Exploración}%
}{%
  \end{alertblock}%
}

\newenvironment{motivacion}{%
  \begin{exampleblock}{Motivación: Ciencia de Datos en la Práctica}%
}{%
  \end{exampleblock}%
}

\newenvironment{retoclase}{%
  \begin{block}{Reto Rápido en Equipo}%
}{%
  \end{block}%
}


\title[MLE y Score]{Sección 06.03: Estimación por Máxima Verosimilitud (MLE) y Score}
\subtitle{La Función de Score y la Normalidad Asintótica del MLE}
\author[J. Castillo Colmenares]{Juliho Castillo Colmenares}
\institute{Tecnológico de Monterrey}
\date{\vspace{-1.2cm}}

\begin{document}

% Slide 1: Portada
\begin{frame}[plain]
  \titlepage
\end{frame}

% Slide 2: Hoja de Ruta
\begin{frame}{Hoja de Ruta --- Capítulo 06: Estimación y su Relación con Ciencia de Datos}
  \footnotesize
  ¿Dónde nos encontramos en el desarrollo modular del capítulo? \pause
  \begin{itemize}\itemsep=0.08em
    \item \textbf{06.01} Estimación Puntual, Insesgadez, Eficiencia y Consistencia \pause
    \item \textbf{06.02} Método de Momentos (MoM) \pause
    \item \textbf{06.03} Estimación por Máxima Verosimilitud (MLE) y Score \emph{(Hoy)} \pause
    \item \textbf{06.04} Intervalos de Confianza para Medias Poblacionales ($Z$ y $t$) \pause
    \item \textbf{06.05} Intervalos de Confianza para Varianzas ($\chi^2$) y Proporciones
  \end{itemize}
  \vspace{-0.05cm}
  \begin{block}{Objetivo de la Sesión}
    Formalizar la función de score y sus propiedades (media cero, identidad de la información), establecer la normalidad asintótica del MLE, y conectar este resultado con la Cota de Cramér-Rao de la Sección 06.01.
  \end{block}
\end{frame}

% Slide 3: Motivación
\begin{frame}{Del Método de Momentos al Estimador Óptimo}
  \footnotesize
  \begin{motivacion}
    En la Sección 06.02 vimos que el MoM es simple pero generalmente subóptimo. El MLE, en cambio, alcanza la eficiencia asintótica óptima. Hoy formalizamos \emph{por qué}: a través de la función de score y su varianza, la Información de Fisher.
  \end{motivacion}

  \vspace{0.15cm}
  \pause
  \begin{itemize}\itemsep=0.1cm
    \item \textbf{Score:} la derivada de la log-verosimilitud; su raíz define al MLE. \pause
    \item \textbf{Identidad de la información:} conecta la varianza del score con la curvatura de la log-verosimilitud. \pause
    \item \textbf{Normalidad asintótica:} el MLE se distribuye aproximadamente Normal para $n$ grande, con varianza igual a la Cota de Cramér-Rao.
  \end{itemize}
\end{frame}

% Slide 4: Función de Score
\begin{frame}{La Función de Score}
  \footnotesize
  \begin{block}{Definición}
    $$U(\theta;X) = \frac{\partial}{\partial\theta}\ln f(X\mid\theta), \qquad U_n(\theta) = \sum_{i=1}^n U(\theta;X_i) = \frac{d}{d\theta}\ell(\theta)$$
    El MLE resuelve la \emph{ecuación de score}: $U_n(\hat\theta_{MLE})=0$.
  \end{block}

  \vspace{0.1cm}
  \pause
  \begin{alertblock}{Propiedades Fundamentales}
    \begin{itemize}\itemsep=0.05cm
      \item \textbf{Media cero:} $E_\theta[U(\theta;X)]=0$.
      \item \textbf{Identidad de la información:} $\Var_\theta[U(\theta;X)] = I(\theta)$.
    \end{itemize}
  \end{alertblock}
\end{frame}

% Slide 5: Normalidad Asintótica
\begin{frame}{Normalidad Asintótica del MLE}
  \footnotesize
  \begin{block}{Teorema}
    $$\sqrt{n}(\hat\theta_{MLE}-\theta) \xrightarrow{d} N\!\left(0,\, \frac{1}{I(\theta)}\right) \implies \hat\theta_{MLE} \overset{\text{aprox.}}{\sim} N\!\left(\theta,\, \frac{1}{nI(\theta)}\right)$$
  \end{block}

  \vspace{0.1cm}
  \pause
  \begin{alertblock}{Conexión con Cramér-Rao}
    La varianza asintótica del MLE coincide exactamente con la Cota Inferior de Cramér-Rao (Sección 06.01): el MLE es \textbf{asintóticamente óptimo}, aunque no sea insesgado ni óptimo para $n$ finito (recuerde $\hat\sigma^2_{MLE}$, sesgado por un factor $(n-1)/n$).
  \end{alertblock}
\end{frame}

% Slide 6: Aplicación Práctica
\begin{frame}{Aplicación: Intervalos de Confianza Aproximados}
  \footnotesize
  La normalidad asintótica permite construir intervalos de confianza sin conocer la distribución exacta de $\hat\theta_{MLE}$: \pause

  \vspace{0.1cm}
  \begin{block}{Fórmula General}
    $$\hat\theta_{MLE} \pm z_{\alpha/2}\sqrt{\dfrac{1}{n I(\hat\theta_{MLE})}}$$
  \end{block}

  \vspace{0.1cm}
  \pause
  \begin{alertblock}{Ejemplo: Poisson}
    Con $I(\lambda)=1/\lambda$: para $n=100$, $\hat\lambda_{MLE}=3.5$, el IC del $95\%$ es $3.5 \pm 1.96\sqrt{3.5/100} \approx [3.13, 3.87]$.
  \end{alertblock}
\end{frame}

% Slide 7: Aplicaciones en Ciencia de Datos
\begin{frame}{Aplicaciones en Ciencia de Datos}
  \footnotesize
  \begin{itemize}\itemsep=0.1cm
    \item \textbf{Regresión logística y redes neuronales:} el entrenamiento vía descenso de gradiente maximiza directamente $\ell(\theta)$; el gradiente \emph{es} el score. \pause
    \item \textbf{Errores estándar de modelos:} el inverso de la matriz de información de Fisher da los errores estándar reportados por software estadístico (\texttt{statsmodels}, \texttt{R}). \pause
    \item \textbf{Pruebas de score (LM):} pruebas de hipótesis basadas únicamente en $U_n(\theta_0)$, sin necesidad de ajustar el modelo alternativo completo. \pause
    \item \textbf{Método delta:} propaga la varianza asintótica del MLE a transformaciones no lineales de los parámetros.
  \end{itemize}
\end{frame}

% Slide 8: Lab Python (Bloque 1)
\begin{frame}[fragile]{Laboratorio en Python: Propiedades del Score}
  \scriptsize
  Verificación de $E[U]=0$ y la identidad de la información para la Exponencial (\texttt{06.03\_maximum\_likelihood\_estimation.py}):
  {\lstset{basicstyle=\fontsize{4pt}{4.8pt}\selectfont\ttfamily,aboveskip=0.1em,belowskip=0.1em}
  \lstinputlisting[firstline=17, lastline=29]{../../code/06_estimacion_estadistica/06.03_maximum_likelihood_estimation.py}}
\end{frame}

% Slide 9: Lab Python (Bloque 2)
\begin{frame}[fragile]{Laboratorio en Python: Normalidad Asintótica}
  \scriptsize
  Verificación Monte Carlo de la normalidad asintótica del MLE de Poisson y construcción de un IC aproximado:
  {\lstset{basicstyle=\fontsize{4pt}{4.8pt}\selectfont\ttfamily,aboveskip=0.1em,belowskip=0.1em}
  \lstinputlisting[firstline=32, lastline=55]{../../code/06_estimacion_estadistica/06.03_maximum_likelihood_estimation.py}}
\end{frame}

% Slide 10: Lab Python (Bloque 3)
\begin{frame}[fragile]{Laboratorio en Python: MLE de Rayleigh y Método Delta}
  \scriptsize
  Derivación del MLE de Rayleigh y verificación del método delta para $\hat\sigma_{MLE}$:
  {\lstset{basicstyle=\fontsize{4pt}{4.8pt}\selectfont\ttfamily,aboveskip=0.1em,belowskip=0.1em}
  \lstinputlisting[firstline=58, lastline=85]{../../code/06_estimacion_estadistica/06.03_maximum_likelihood_estimation.py}}
\end{frame}

% Slide 11: Salida Terminal
\begin{frame}[fragile]{Laboratorio en Python: Salida en Terminal}
  \scriptsize
  Salida estándar generada al ejecutar el script del laboratorio en Python:
  \vspace{0.02cm}
  \begin{block}{Salida Estándar en Consola}
    \fontsize{4pt}{5pt}\selectfont
    \begin{verbatim}
=== Block 1: Score Function Properties (Exponential) ===
E[U]=0.00094 (theo 0) | Var[U]=0.249 (theo I(lambda)=0.25)

=== Block 2: Asymptotic Normality (Poisson) ===
Var[lambda_MLE]=0.0350 (theo 1/(nI)=0.0350)
Approx. 95% CI (n=100, lambda_hat=3.5): [3.133, 3.867]

=== Block 3: Rayleigh MLE and Delta Method ===
Var(sigma_MLE)=0.0200 (theo sigma^2/(4n)=0.0200)
Delta-method Var(sigma) from Var(sigma^2): 0.0200
    \end{verbatim}
  \end{block}
\end{frame}

% Slide 12A: Ejercicio Nivel 1 (Enunciado)

% Slide 12B: Ejercicio Nivel 1 (Resolución)

% Slide 13A: Ejercicio Nivel 2 (Enunciado)

% Slide 13B: Ejercicio Nivel 2 (Resolución)

% Slide 14A: Ejercicio Nivel 3 (Enunciado)

% Slide 16: Cierre
\begin{frame}{Cierre de Sección: MLE y Score}
  \begin{reflexion}
    La función de score es el hilo conductor que conecta la construcción del MLE (ecuación de score), la Información de Fisher (varianza del score), y la Cota de Cramér-Rao (Sección 06.01). Juntas, estas ideas explican por qué el MLE es, en el límite, el mejor estimador posible.
  \end{reflexion}

  \vspace{0.2cm}
  \pause
  \begin{block}{Lecciones Clave}
    \begin{itemize}\itemsep=0.08cm
      \item \textbf{Score:} $U(\theta;X)=\partial\ln f/\partial\theta$; media cero, varianza $I(\theta)$.
      \item \textbf{Normalidad asintótica:} $\hat\theta_{MLE} \overset{\text{aprox.}}{\sim} N(\theta, 1/(nI(\theta)))$ --- la varianza óptima de Cramér-Rao.
      \item \textbf{Aplicación:} intervalos de confianza aproximados y el método delta para transformaciones.
    \end{itemize}
  \end{block}
\end{frame}

% Slide 17: Perspectiva Modular
\begin{frame}{Perspectiva Modular --- El Próximo Reto}
  \scriptsize
  \begin{retoclase}
    Con MoM y MLE dominados, la Sección 06.04 aplicará esta maquinaria a un problema muy concreto: construir intervalos de confianza para la media poblacional $\mu$, usando $Z$ cuando $\sigma$ es conocida y $t$ cuando no lo es --- retomando directamente la distribución $t$ de Student del Capítulo 05.
  \end{retoclase}

  \vspace{0.08cm}
  \pause
  \begin{alertblock}{Preparación Recomendada}
    \begin{itemize}\itemsep=0.06cm
      \item Repasa el intervalo de confianza $t$ de la Sección 05.04.
      \item Reflexiona sobre cómo la normalidad asintótica del MLE generaliza esa construcción a \emph{cualquier} parámetro, no solo a $\mu$.
    \end{itemize}
  \end{alertblock}
\end{frame}

\end{document}
